\documentclass{situationist-preprint}

\title{The Expanding Circle Revisited}
\msid{MTE-19-b}
\msfamily{MTE-19}
\msdiscipline{Metaethics}
\mstype{Companion manuscript}
\mskeywords{moral progress; expanding circle; Singer; Lecky; inclusivity; fixed standard; regress}

\begin{document}
\maketitle

\begin{abstract}
The image of the expanding moral circle, due to Lecky and developed by Singer, is the natural model of moral progress, and MTE-19-a makes it precise as a measure against a fixed standard. We revisit the image in the light of three difficulties: that circles can expand in the wrong direction, that much progress concerns treatment within the circle rather than its boundary, and that expansion has proved fragile. Each difficulty is met by a feature of the precise account: the measure counts only beings in the correct circle, the graded measure counts improvements within it, and the fixity of the standard is what allows contraction to be recognized as regress.
\end{abstract}

\section{The image and its precise form}

\citet{Lecky1869} described the history of morals as a widening of the circle of beings whose interests are counted, from family to tribe to nation to humanity and, at length, to animals. \citet{Singer1981} gave the image a mechanism: reasoning about conduct, once begun within a group, requires justification in impartial terms, and impartiality, once accepted, cannot be confined to the group in which it was first practiced. MTE-19-a turns the image into a measure. The moral circle of a society is the set of beings to whom its effective norms give due consideration (Definition~2.1 of MTE-19-a); the standard is the correct circle (Definition~2.2), fixed independently of any society's norms (Assumption~2.3); and progress is inclusion of more of the standard (Definitions~2.4 and~2.5). Theorem~3.5 of MTE-19-a shows that progress so defined is real.

The image has been criticized on three grounds. We take them in turn.

\section{Three difficulties}

\paragraph{Expansion in the wrong direction.} A society may extend its consideration to beings that are not owed it, such as idols or the spirits of ancestors, and may do so at the expense of beings that are. A count of the beings a society considers would register such a change as expansion. The measure of MTE-19-a does not, since it counts only the members of the circle that are also in the standard (Definition~2.4 of MTE-19-a). By ETH-07, the standard is the class of sentient beings, with weight proportional to their capacity for welfare. Consideration given to beings outside it does not raise the measure, and if it is given at the cost of beings within it, as when resources owed to the living are diverted to the dead, the measure falls. The image of expansion was always shorthand for expansion toward the correct boundary, and the precise account says so.

\paragraph{Progress within the circle.} Much of what we count as moral progress is not the admission of new beings but better treatment of those already counted: the abolition of judicial torture, of debtors' prisons, of the subordination of wives to husbands. The all-or-nothing measure misses these. The graded measure of Remark~2.6 of MTE-19-a does not: it assigns each being the proportion of its due weight that the society's norms give it, and improvement within the circle raises the proportion. Singer's mechanism explains both kinds of change in the same way, since the demand for impartial justification bears on how a counted interest is weighed as much as on whose interests are counted.

\paragraph{Fragility.} \citet{BuchananPowell2018} argue that our evolved moral psychology is tuned to the in-group and that inclusive norms emerge, and persist, only under favorable conditions of security and prosperity; when those conditions fail, the circle contracts. The historical record bears them out, and MTE-19-a records two cases of contraction (Proposition~4.2 and the discussion following it). This is a difficulty for the image only if it is read as a law of history. It is not one for the measure. On the contrary, contraction can be recognized as regress only because the standard is fixed: if the standard were set by the latest stage of a history, no contraction at the end of it could count as regress, by definition (Proposition~3.4 of MTE-19-a). Fragility is therefore evidence for the fixed standard, not against the reality of progress.

\section{Objection and reply}

\begin{objection}
If the standard is the correct circle, the measure can be applied only by someone who knows what the correct circle is. Since we do not know it with certainty, we cannot know whether we have progressed.
\end{objection}

\begin{reply}
The measure requires knowledge of the standard only for the beings in the class relevant to a comparison, and our knowledge there is often secure. That the enslaved were owed consideration is as certain as any moral claim, and it suffices to establish that abolition was progress. Where our knowledge is less secure, as about the capacities for welfare of some animals, our judgments of progress are correspondingly uncertain. This is the ordinary situation of measurement against a standard known imperfectly, and it is the same in kind as the fallible but reliable knowledge of moral truths that MTE-08 defends.
\end{reply}

\section{Conclusion}

The expanding circle survives its revision. Expansion counts as progress only when it is expansion toward the correct circle; progress includes improvements of treatment within the circle; and the fragility of inclusion is real, but it is precisely the fixity of the standard that allows contraction to be identified as regress. In this, moral progress differs from what can be said of artistic value, for which, by AES-16, there is progress in means but none in value, since there is no fixed standard of the same kind. Moral progress is real and measurable, and neither steady nor irreversible.

\begin{thebibliography}{9}\small
\bibitem[Buchanan and Powell(2018)]{BuchananPowell2018} Buchanan, A. and Powell, R. (2018). \emph{The Evolution of Moral Progress: A Biocultural Theory}. Oxford University Press.
\bibitem[Lecky(1869)]{Lecky1869} Lecky, W.~E.~H. (1869). \emph{History of European Morals from Augustus to Charlemagne}. Longmans, Green, and Co.
\bibitem[Singer(1981)]{Singer1981} Singer, P. (1981). \emph{The Expanding Circle: Ethics and Sociobiology}. Farrar, Straus and Giroux.
\end{thebibliography}

\end{document}
